\begin{textsample}{NEG Dim 4 – 2016\_c5 – Score 2.00 – 2016\_c5/enade\_2016\_tec\_rad\_371c7e6d01.txt}  \label{ex:f4_neg_256}
CAPÍTULO 5 – DISTRIBUIÇÃO DOS CONCEITOS 5.1 - PANORAMA NACIONAL DA DISTRIBUIÇÃO DOS CONCEITOS O Conceito Enade é um indicador de qualidade que avalia o desempenho dos estudantes a partir dos resultados do Enade.
Os valores possíveis vão de 1 ( pior situação ) a 5 ( melhor situação ).
As Unidades de Observação com apenas um ou sem nenhum concluinte participante não obtêm o Conceito Enade, ficando Sem Conceito ( SC ).
A Tabela 5.1 apresenta a quantidade e distribuição de cursos de Tecnologia em Radiologia participantes do Enade/2016, por faixa de conceito e Grande Região.
Enfatiza -se, mais uma vez, que a diferença entre os cursos tabulados neste capítulo e no capítulo 2 corresponde aos cursos sem conceito, ou seja àqueles, em princípio, sem alunos concluintes que houvessem participado da prova ou àqueles que tivessem tão somente um aluno concluinte19.
Observando -se os dados da Tabela 5.1, nota -se que, dos 95 cursos participantes, 34 ( 35,8 \% ) classificaram -se com conceito 2, o valor modal.
Esse foi também o conceito modal em quatro regiões : Norte ( 50,0 \% ),, Sul ( 38,5 \% ), Centro-Oeste ( 50,0 \% ) e Nordeste ( 25,0 \% ) com distribuição multimodal nos conceitos 1, 2 e 4 ( 25,0 \% ).
Na região Sudeste, a moda foi o conceito 3, atribuído para 38,6 \% dos cursos da região.
O conceito 3 foi o segundo mais frequente em nível nacional ( 29,5 \%, correspondendo a 28 cursos ) e o conceito 4, o terceiro ( 16,8 \%, correspondendo a 16 cursos ).
Houve, ainda, seis cursos ( 6,3 \% ) que receberam conceito 5 e nove cursos ( 9,5 \% ) que receberam conceito 1.
Dos 95 cursos de Tecnologia em Radiologia, dois cursos ficaram sem conceito ( SC ). 19 Estes cursos com somente um concluinte não foram considerados no capítulo 2.
A região Norte participou com oito cursos ou 8,4 \% do total nacional.
Desses, quatro cursos receberam o conceito 2, o que equivale a 50,0 \% do total regional, além de ser o conceito modal na região, como já comentado.
A um curso ( 12,5 \% ) atribuiu -se conceito 1 e a dois cursos ( 25,0 \% ), conceito 3.
Um curso ( 12,5 \% ) ficou sem conceito ( SC ).
Nenhum curso ficou com conceito 4 ou 5 nessa região.
A região Nordeste participou com 20 cursos ou 21,1 \% do total nacional.
Nota -se que tais cursos apresentaram distribuição multimodal, uma vez que as maiores frequências se concentraram em três conceitos, os conceitos 1, 2 e 4, atribuídos, cada qual, a cinco cursos ( 25,0 \% cada ).
Os conceitos restantes 3 e 5 foram atribuídos, respectivamente, a quatro e a um curso ( 20,0 \% e 5,0 \%, respectivamente ).
Nenhum curso ficou sem conceito ( SC ).
Dos 44 cursos participantes da região Sudeste ( 46,3 \% do total nacional ), o conceito 3 foi o conceito modal, como já anteriormente assinalado, tendo sido obtido por 17 cursos ( 38,6 \% ).
O conceito 1 foi atribuído a dois cursos ( 4,5 \% ) e o conceito 2, a 15 cursos ( 34,1 \% ).
Receberam os conceitos 4 e 5, respectivamente, cinco cursos ( 11,4 \% ) e quatro cursos ( 9,1 \% ).
Nessa região, um curso ( 2,3 \% ) ficou sem conceito ( SC ).
Os 13 cursos da região Sul corresponderam a 13,7 \% do total nacional.
Como já indicado, houve predominância do conceito 2, atribuído a 38,5 \% dos cursos da região ( conceito modal ), ou seja, atribuído a cinco dos 13 cursos participantes na região Sul.
O conceito 3 foi atribuído a quatro cursos ( 30,8 \% ), o conceito 4, a três cursos ( 23,1 \% ) e o conceito 5, a um curso ( 7,7 \% ).
Nenhum curso recebeu conceito 1 ou ficou sem conceito ( SC ) na região Sul.
Já dos dez cursos participantes na região Centro-Oeste ( 10,5 \% do total nacional ), cinco receberam conceito 2, o conceito modal ( 50,0 \% ), como já fora destacado.
A dois cursos foram atribuídos conceitos 1 e 3 ( um curso, cada ), e a outros Nessa região, nenhum curso ficou sem conceito ( SC ) ou com conceito 5.
O Gráfico 5.1 apresenta a distribuição do Conceito Enade segundo Grande Região : poligonais mais à esquerda representam Grandes Regiões com pior distribuição, e poligonais mais à direita, Grandes Regiões com melhores conceitos.
A região Norte ( linha azul ) apresenta os cursos com conceitos mais baixos, ao passo que as regiões Sudeste ( linha vermelha ) e Sul ( linha roxa ) apresentam as distribuições com valores maiores.
Cursos em IES situadas nas regiões Nordeste ( linha verde ) e Centro-Oeste ( linha preta ) apresentam uma situação intermediária. 5.2 - CONCEITOS POR CATEGORIA ADMINISTRATIVA, POR MODALIDADE DE ENSINO E GRANDE REGIÃO A Tabela 5.2 apresenta a distribuição dos cursos participantes do Enade/2016 de Tecnologia em Radiologia, por Categoria Administrativa e Modalidade de Ensino, de acordo com os conceitos por eles alcançados, segundo as Grandes Regiões brasileiras.
Dos 95 cursos participantes, dez ( 10,5 \% ) eram ministrados em Instituições Públicas e 85 ( 89,5 \% ), em Instituições Privadas.
De acordo com as informações da Tabela 5.2, em termos nacionais, entre Instituições Públicas e Privadas, observa -se que, dos seis cursos avaliados com conceito 5, dois eram oferecidos em IES Públicas e quatro, em IES Privadas.
Dos dez cursos participantes de IES Públicas, o conceito 4 foi o valor modal, atribuído a quatro cursos ( 40,0 \% ).
Entre os demais cursos de Instituições Públicas participantes, um curso ( 10,0 \% da categoria ) obteve conceito 2, três cursos ( 30,0 \% ) receberam conceito 3 e, como já comentado, dois cursos ( 20,0 \% ) receberam conceito 5.
Nenhum curso ficou sem conceito ( SC ).
Na rede privada, o conceito modal foi 2, atribuído a 33 cursos, o correspondente a 38,8 \% dos 85 cursos da categoria.
Entre os demais cursos participantes, nove ( 10,6 \% ) receberam conceito 1, e 25 ( 29,4 \% ), conceito 3.
O conceito 4 foi atribuído a 12 cursos ( 14,1 \% ) e o conceito 5, a quatro cursos ( 4,7 \% ).
Dois cursos ( 2,4 \% ) ficaram sem conceito ( SC ).
Os dados referentes à distribuição dos cursos segundo a Modalidade de Ensino revelam que apenas um dos 95 cursos de Tecnologia em Radiologia participantes do Enade 2016 foi oferecido na modalidade Educação a Distância, equivalente 1,1 \% dos cursos.
A esse curso, foi atribuído conceito 4.
A distribuição dos conceitos nos cursos presenciais foi quase idêntica a do Brasil como um todo, considerando -se que 98,9 \% dos cursos eram desta modalidade.
Na análise por região, observa -se que, na região Norte, as Instituições Privadas participaram com oito cursos ( 100,0 \% do total regional ), dos quais um ( 12,5 \% ) ficou sem conceito ( SC ).
Nenhum curso recebeu conceito 4 ou 5.
O conceito modal para as Instituições Privadas na região foi 2, atribuído a quatro cursos, correspondendo a 50,0 \% dessa categoria na região.
Um único curso ( 12,5 \% ) foi avaliado com conceito 1 e dois cursos receberam conceito 3 ( 25,0 \% ).
As Instituições Públicas não tiveram participação na região Norte.
Na região Nordeste, a rede privada concentrou 16 dos 20 cursos participantes, o equivalente a 80,0 \% do total da região, a menor das proporções desta categoria entre as regiões.
Dentre os cursos oferecidos por IES Privadas no Nordeste, nota -se que o conceito modal foi 1, atribuído a cinco cursos ( 31,3 \% ).
Oito cursos obtiveram conceitos 2 e 3 ( quatro cursos, cada ) e outros três, obtiveram conceito 4.
Nenhum curso ficou sem conceito ( SC ) ou conceito 5.
As Instituições Públicas da região Nordeste participaram com quatro cursos ( 20,0 \% do total da região ), a maior das proporções desta categoria nas regiões.
Desses, dois obtiveram conceito 4.
Os demais foram avaliados com conceitos 2 e 5 ( um curso, cada ).
Nenhum curso ficou sem conceito ( SC ) ou com conceito 1 ou 3.
Na região Sudeste, a proporção de cursos da rede privada foi de 90,9 \%.
Nessa categoria e região, o conceito modal foi 2, atribuído a 15 cursos.
Um curso ficou sem conceito ( SC ).
Os demais foram avaliados com conceito 1 ( dois cursos ), conceito 3 ( 14 cursos ) e conceitos 4 e 5 ( quatro cursos, cada ).
Entre os quatro cursos oferecidos por Instituições Públicas na região Sudeste, três receberam conceito 3, o conceito modal, e um ficou com conceito 4.
Na região Sudeste, nesta categoria, nenhum curso ficou sem conceito ( SC ) ou com os demais conceitos.
As Instituições Privadas concentraram 11 dos 13 cursos participantes da região Sul, 84,6 \% do total regional.
Desses, cinco obtiveram conceito 2, o conceito modal.
Nesta combinação de Categoria Administrativa e Grande Região, os conceitos 3 e 4 foram atribuídos a quatro e a dois cursos, respectivamente.
Nenhum curso ficou sem conceito ( SC ) ou recebeu conceito 1 ou 5.
As Instituições Públicas na região Sul participaram com somente dois cursos, um foi avaliado com conceito 4 e outro com conceito 5.
Na região Centro-Oeste, todos os dez cursos participantes eram de Instituições Privadas ( 100,0 \% em termos regionais ).
Desses, cinco receberam conceito 2, o conceito modal.
Entre os demais, dois receberam conceitos 1 e 3 ( um curso, cada ) e três receberam conceito 4.
Nenhum curso ficou sem conceito ( SC ) ou recebeu conceito 5.
As Instituições Públicas não tiveram participação na região Centro-Oeste.
Como já comentado apenas um curso foi oferecido na modalidade a distância.
Este curso estava situado na região Sudeste, correspondendo a 2,3 \% dos cursos na região.
A esse curso, mantido por IES situada na região foi atribuído conceito 4.

% matched lemmas: 
\end{textsample}
